% Canales 1 a 11 de la banda de 2,4 GHz (fc = 2407 + 5n MHz, ancho aprox. 22 MHz)
% Eje x: MHz por encima de 2400
\begin{tikzpicture}[x=0.15cm, y=1cm]
  % Canales superpuestos
  \foreach \n in {2,3,4,5,7,8,9,10}{
    \pgfmathsetmacro{\fc}{7+5*\n}
    \draw[gray!65, dashed] plot[domain=-11:11, samples=41] ({\fc+\x},{1.8*cos(\x*180/22)});
    \node[font=\scriptsize, text=gray!70!black, above] at (\fc,1.8) {\n};
  }
  % Canales sin solapamiento
  \foreach \n in {1,6,11}{
    \pgfmathsetmacro{\fc}{7+5*\n}
    \fill[FigAcento, opacity=0.18] plot[domain=-11:11, samples=41] ({\fc+\x},{1.8*cos(\x*180/22)}) -- cycle;
    \draw[FigAcento, very thick] plot[domain=-11:11, samples=41] ({\fc+\x},{1.8*cos(\x*180/22)});
    \node[font=\scriptsize\bfseries, text=FigAcento, above] at (\fc,1.8) {\n};
  }
  % Eje de frecuencia
  \draw[-{Stealth[length=2mm]}] (-2,0) -- (90,0);
  \foreach \f/\etiqueta in {0/2400, 12/2412, 37/2437, 62/2462, 83.5/{2483,5}}{
    \draw (\f,0) -- (\f,-0.12) node[below, font=\scriptsize] {\etiqueta};
  }
  \node[font=\scriptsize, anchor=north east] at (90,-0.5) {Frecuencia (MHz)};
  % Leyenda
  \fill[FigAcento, opacity=0.18] (2,-1.3) rectangle (6,-1.0);
  \draw[FigAcento, very thick] (2,-1.3) rectangle (6,-1.0);
  \node[font=\scriptsize, anchor=west] at (6.5,-1.15) {Canales sin solapamiento (1, 6 y 11)};
  \draw[gray!65, dashed] (42,-1.15) -- (48,-1.15);
  \node[font=\scriptsize, anchor=west] at (48.5,-1.15) {Canales que se superponen};
\end{tikzpicture}
